\documentclass[11pt,letterpaper]{article}
\usepackage[T1]{fontenc}
\usepackage{lmodern}
\usepackage[left=.95in,right=.95in,top=.89in,bottom=.91in]{geometry}
\usepackage{amsmath,amssymb,amsthm,mathtools,microtype}
\usepackage{xcolor}
\usepackage[colorlinks=true,linkcolor=blue!40!black,citecolor=blue!40!black,
  urlcolor=blue!40!black]{hyperref}
\hypersetup{pdftitle={The alternating sum of reciprocal Fibonacci cubes is irrational},
  pdfauthor={twinheterodox},pdfsubject={Number theory}}
\newtheorem{theorem}{Theorem}[section]
\newtheorem{corollary}[theorem]{Corollary}
\newtheorem{proposition}[theorem]{Proposition}
\newtheorem{lemma}[theorem]{Lemma}
\numberwithin{equation}{section}
\newcommand{\Q}{\mathbb Q}
\newcommand{\Z}{\mathbb Z}
\newcommand{\Ord}{\mathcal O}
\DeclareMathOperator{\Nm}{N}
\title{\fontsize{16}{20}\selectfont\mbox{The alternating sum of reciprocal Fibonacci cubes is irrational}}
\author{twinheterodox}
\date{18 November 2026}

\begin{document}
\maketitle

\begin{abstract}
We prove that the alternating sum of the reciprocal cubes of the Fibonacci numbers
is irrational. The result extends to the Lucas sequences of the first kind
associated with real quadratic units greater than one, including the Pell
numbers. For every such sequence and every fixed positive integer \(k\), the
reciprocal cubes with indices \(k,2k,3k,\ldots\), taken with signs
\(+,-,+,\ldots\), also have an irrational sum.
\end{abstract}

\section{The alternating sum}

Let \(F_1=F_2=1\) and \(F_{n+2}=F_{n+1}+F_n\). Write
\[
 A_F^+=\sum_{n\ge1}\frac{1}{F_n^3},\qquad
 A_F^-=\sum_{n\ge1}\frac{(-1)^{n-1}}{F_n^3}.
\]
Absolute convergence gives \(A_F^++A_F^-=2\sum_{n\ge1}F_{2n-1}^{-3}\).
Duverney, Ke.~Nishioka, Ku.~Nishioka and Shiokawa proved that the
odd-index sum is transcendental \cite{DNNS}\cite[Theorem~4, p.~95]{Kyoto}.
Thus at least one of \(A_F^+\) and \(A_F^-\) is transcendental.
Krisman proved \(A_F^+\) irrational \cite{KFib} and extended that
result to ordinary reciprocal cube sums of Lucas sequences with a
positive dominant unit \cite[Theorem~1.1]{KLucas}. These results
leave open whether \(A_F^-\) is rational \cite[p.~3]{KLucas}.

For integers \(P,Q\), let \(U_n=U_n(P,Q)\) be the Lucas sequence of the
first kind,
\[
 U_0=0,\qquad U_1=1,\qquad U_{n+2}=P U_{n+1}-Q U_n.
\]
\begin{theorem}\label{thm:main}
Suppose that \(Q=-1\) and \(P\ge1\), or that \(Q=1\) and \(P\ge3\).
Then
\[
                 A^-_{P,Q}=\sum_{n\ge1}\frac{(-1)^{n-1}}{U_n(P,Q)^3}
                 \quad\text{is irrational}.
\]
\end{theorem}
\noindent
Changing the sign of the trace gives, directly from the recurrence,
\[
   U_n(-P,Q)=(-1)^{n-1}U_n(P,Q)\qquad(n\ge1).
\]
Cubing preserves this sign. Thus \(A^-_{P,Q}\) is the ordinary
reciprocal cube sum for trace \(-P\). In the Fibonacci case
\((P,Q)=(1,-1)\), the negative-trace sequence begins \(1,-1,2,-3,5,\ldots\).
Its dominant root is negative, explaining the signed reciprocal root
in Section~\ref{sec:moments}. The pair \((2,-1)\) gives the Pell case.
For \(Q=-1\), Elsner, Shimomura and Shiokawa treat
alternating reciprocal \emph{even} powers of first- and second-kind
Lucas sequences \cite[Theorems~2 and~4]{ESS}. Here the power is three.

\begin{corollary}\label{cor:stride}
Under the hypotheses of Theorem~\ref{thm:main}, for each integer \(k\ge1\),
\[
                  \sum_{n\ge1}\frac{(-1)^{n-1}}{U_{kn}(P,Q)^3}
                  \quad\text{is irrational}.
\]
\end{corollary}
The determinant method produces polynomials in a single unknown value.
Under a rationality assumption, clearing their denominators gives
algebraic integers in a real quadratic field. A nonzero algebraic
integer has absolute norm at least one, so estimates at \emph{both}
real embeddings can yield a contradiction. Zudilin used polynomial
Hankel determinants in an irrationality criterion
\cite[Proposition~2]{Zudilin}. Uvarov related the orthogonal polynomials
before and after a rational change of weight \cite[\S1]{Uvarov}.
Krattenthaler gave a determinant identity for the modified moments
\cite[Theorem~1]{Kratt}.
The particular moving denominator and interpolation argument used here
come from the \(q\)-zeta construction \cite[Sections~2 and~5]{QZeta}
and Krisman's adaptation of it to reciprocal cubes \cite{KFib,KLucas}.

The alternating sign amounts to changing the sign of the small unit in
Krisman's construction. His cyclotomic denominator is an identity in a
formal variable, so it still clears every coefficient. The main analytic
point is at the conjugate embedding. The conjugate pole values turn out
to be evaluations of a \emph{single} holomorphic function formed from the
original signed moments. A contour integral combines them before their
sizes are estimated. For a determinant of size \(K\), measure logarithms
in units of \(K^3\log\eta\), where \(\eta>1\) is the dominant unit.
The first embedding saves \(2/9\), the conjugate embedding saves \(1\),
and clearing the denominators costs at most \(4393/3675\). The margin is
\begin{equation}\label{eq:margin}
                    \frac29+1-\frac{4393}{3675}=\frac{296}{11025}.
\end{equation}
We prove the estimates for the actual signed value and show that
infinitely many of the corresponding determinants are nonzero.

\section{The signed unit and its moments}\label{sec:moments}

Put \(D=P^2-4Q\), \(\eta=(P+\sqrt D)/2\), and \(u=\eta^{-1}\).
The two roots of \(x^2-Px+Q\) are \(\eta>1\) and \(Q/\eta\).
They are algebraic units. A rational root would be an integer
dividing \(Q=\pm1\), so the field \(E=\Q(\sqrt D)\) is quadratic.
Conversely, the trace of a real quadratic unit \(\eta>1\) of norm
\(Q\) is the integer \(\eta+Q/\eta\), which is positive for \(Q=-1\)
and exceeds two for \(Q=1\). This gives exactly the parameter ranges
in Theorem~\ref{thm:main}. Let \(\sigma\) denote the nontrivial
automorphism of \(E\) and put
\[
     \tau=-u,\qquad w=Q\tau^2,\qquad \rho=Q\tau.
\]
Here \(|w|<1\), \(|\rho|<1\), \(w=\rho\tau\), and \(\rho w=\tau^3\).
Changing \(\tau\) from \(u\) to \(-u\) fixes the nodes \(w^n\) and
multiplies the weights \(\rho^n\) by \((-1)^n\). Binet's
formula gives \(U_n=\eta^n(1-w^n)/\sqrt D\). Thus all the \(U_n\)
with \(n\ge1\) are positive and
\begin{equation}\label{eq:period}
 \psi_N=\frac{\tau^{3N}}{(1-w^N)^3},\qquad
 S_a=\sum_{N=1}^{a}\psi_N,\qquad \Xi=\sum_{N\ge1}\psi_N,\qquad
 A^-_{P,Q}=-D\sqrt D\,\Xi.
\end{equation}
Here \(\tau^{3N}=(-1)^N\eta^{-3N}\), which accounts for the minus sign
in the formula for \(A^-_{P,Q}\). The series is absolutely convergent.
The finite quantities
satisfy
\begin{equation}\label{eq:conj}
 \sigma\tau=Q/\tau,\qquad \sigma w=w^{-1},\qquad
 \sigma\rho=\tau^{-1},\qquad \sigma\psi_N=-\psi_N .
\end{equation}
For the last identity, factor \(-w^{-N}\) out of \(1-w^{-N}\).
We shall use \(\sigma\Xi=-\Xi\) only under the hypothesis
\(A^-_{P,Q}=v/m\), \(v\in\Z,\ m\ge1\). In that case
\begin{equation}\label{eq:rational}
        mD^2\Xi=-v\sqrt D\in\Ord=\Z[\eta],
        \qquad \sigma\Xi=-\Xi ,
\end{equation}
since \(\sqrt D=2\eta-P\). This argument uses the value
of the convergent series as a whole.

We want polynomial moments in \(E\) and simple-pole moments that recover
tails of \(\Xi\). The Euler operator \(\theta=z\,d/dz\) supplies the
cubed denominator in \eqref{eq:period} through the identity
\(\theta(\theta+1)(1-cz)^{-1}=2cz/(1-cz)^3\).
At the node \(w^n\), weighting the numerator by \(\rho^n\) supplies
\((\rho w)^n=\tau^{3n}\). Thus for \(g\) holomorphic on a
neighborhood of \(|z|\le |w|\), set
\[
       L(g)=\sum_{n\ge1}\rho^n[\theta(\theta+1)g](w^n).
\]
The derivatives of \(g\) are bounded on the indicated disk, so the sum
converges absolutely. With \(\alpha_j=\rho w^j\), summing gives
\begin{align}
 L(z^j)&=\nu_j=\frac{j(j+1)\alpha_j}{1-\alpha_j}
               &&(j\ge0),\label{eq:monomials}\\
 L\big((1-w^az)^{-1}\big)&=2\rho^{-a}(\Xi-S_a)
               &&(a\ge1).\label{eq:poles}
\end{align}
Indeed, in the second identity put \(N=n+a\) and use
\(w^a\tau^{-3a}=\rho^{-a}\). This is the signed specialization of
Krisman's moment identities \cite[Section~2]{KLucas}.

For \(K\ge2\) let \(D_K(z)=\prod_{a=1}^K(1-w^a z)\). On rational
functions of the form \(p/D_K\), define \(L_X\) by ordinary partial
fractions. Use \eqref{eq:monomials} on the polynomial quotient and
\eqref{eq:poles} with \(\Xi\) replaced by \(X\) on the simple fractions.
Set
\begin{equation}\label{eq:det}
           \Delta_{K,\tau}(X)
            =\det_{0\le i,j<K} L_X\big(z^{i+j}/D_K(z)\big).
\end{equation}
Each entry is affine in \(X\), hence this polynomial has degree at
most \(K\). At \(X=\Xi\), all its entries are values of the actual
functional \(L\). Replacing \(\tau\) in the finite formulas by an
indeterminate \(T\) defines \(\Delta_{K,T}(X)\in\Q(T)[X]\). This
formal version will let us use the same denominator at both embeddings.

\section{The denominator and the first embedding}\label{sec:first}

Let \(\Phi_d\) denote the \(d\)th cyclotomic polynomial and \(\varphi\)
Euler's totient. The arithmetic input to our proof is
\cite[Proposition~4.4 and the proof of Proposition~4.6]{KLucas}.
For each fixed \(Q=\pm1\) it supplies finitely supported nonnegative
integers \(e_{K,d}\) such that
\begin{align}
 C_K(T)&=\prod_d\Phi_d(T)^{e_{K,d}},&
                 C_K(T)\Delta_{K,T}(X)&\in\Z[T,T^{-1}][X],
                 \label{eq:clearing}\\
 \sum_d e_{K,d}&=O(K^2\log K),&
 \sum_d\varphi(d)e_{K,d}&\le\frac{4393}{3675}K^3+O(K^2).
                 \label{eq:exponents}
\end{align}
The weight \(\varphi(d)\) in \eqref{eq:exponents} is the degree of
\(\Phi_d\). The coefficient \(4393/3675\) comes from adding the
contribution of colliding poles, at most
\(\tfrac{11}{8}\prod_{p\le7}(1-p^{-2})\) with the product over primes,
to the contribution \(1/3\) from monomial moments. Since
\eqref{eq:clearing} holds before evaluating \(T\), we may substitute
the negative unit \(\tau\) and its conjugate. Each evaluation has
coefficients in \(\Ord\). The accompanying programs check
\eqref{eq:clearing} coefficientwise for \(2\le K\le5\) and both
signs of \(Q\) \cite{checks}.

\begin{lemma}\label{lem:normcost}
For each of the units in Theorem~\ref{thm:main},
\[
 \log|C_K(\tau)C_K(\sigma\tau)|
     \le \frac{4393}{3675}K^3\log\eta+O(K^2\log K).
\]
\end{lemma}
\begin{proof}
For \(|x|=u\), the product identity
\(\Phi_d(x)=\prod_{a\mid d}(x^a-1)^{\mu(d/a)}\) gives
\[
 \big|\log|\Phi_d(x)|\big|\le E_u,\qquad
            E_u=\sum_{a\ge1}-\log(1-u^a)<\infty .
\]
Here \(\mu\) is the M\"obius function and
\(\big|\log|1-x^a|\big|\le-\log(1-u^a)\) for both real signs of \(x\).
For \(|x|=\eta\), factor \(x^a\) out of \(x^a-1\). The same bound
applies after subtracting \(\varphi(d)\log\eta\). The error in the
logarithm of each factor's norm is therefore at most \(2E_u\).
Sum with weights \(e_{K,d}\) and use \eqref{eq:exponents}.
\end{proof}

All subsequent error constants may depend on the fixed unit. In
logarithmic upper bounds a zero determinant has logarithm \(-\infty\).
\begin{proposition}\label{prop:first}
At the first real embedding,
\[
             \log|\Delta_{K,\tau}(\Xi)|
                    \le -\frac29 K^3\log\eta+O(K^2).
\]
\end{proposition}
\begin{proof}
The functional uses derivatives through order two. A polynomial with a
triple zero at a node contributes zero there, even after multiplication
by a factor holomorphic at that node. This explains the monic basis
from \cite[Proposition~3.1]{KLucas},
\[
 V_i(z)=\prod_{\ell=1}^i\big(z-w^{\lceil\ell/3\rceil}\big),\qquad
 R_i=\lceil i/3\rceil,\qquad b_i=\sum_{\ell=1}^i\lceil\ell/3\rceil .
\]
Using \(V_i,V_j\) in the rows and columns of \eqref{eq:det} preserves
the determinant. When \(n<\max(R_i,R_j)\), at least one factor \(V_i\)
or \(V_j\) has a triple zero at \(w^n\). At every other node, each
factor \(z-w^b\) at \(z=w^n\) has modulus at most \(2|w|^b\).
Applying \(\theta\) once or twice to that factor gives \(z\), of
modulus at most \(|w|^b\). The reciprocal \(1/D_K\) and its first
two Euler derivatives are bounded uniformly on \(|z|\le|w|\). This
follows from Cauchy's estimates on the unit circle, where
\(|D_K(z)|\ge\prod_{a\ge1}(1-|w|^a)>0\).

The product rule and \(\sum_n|\rho|^n<\infty\) now give
\[
       |L(V_iV_j/D_K)|
        \le C(i+j+1)^2\,2^{i+j}|w|^{b_i+b_j}.
\]
Extract \(|w|^{b_i}\) and \(|w|^{b_j}\) from the row and column
bounds. The remaining factors contribute \(O(K^2)\) to the logarithm
by Hadamard's inequality. As \(b_i=i^2/6+O(i)\), we have
\(2\sum_{i<K} b_i=K^3/9+O(K^2)\). Finally,
\(\log|w|=-2\log\eta\), which gives the asserted coefficient.
\end{proof}

\section{A contour for the conjugate determinant}\label{sec:second}

The conjugate of \(w\) has modulus greater than one. We therefore
apply \(\sigma\) to the finite moment formulas, keeping \(X\) formal,
and use their partial fractions to estimate \(\Delta_{K,\sigma\tau}\).
Denote the resulting functional by \(L_X^\sigma\). The conjugated
denominator vanishes at the points \(r_a=w^a\) inside the unit circle.
Indeed,
\begin{equation}\label{eq:denconj}
 \sigma D_K(z)=(-1)^K w^{-K(K+1)/2}F_K(z),\qquad
 F_K(z)=\prod_{a=1}^K(z-r_a).
\end{equation}
Passing to the reciprocal denominator gives a factor of absolute
value \(|w|^{K(K+1)/2}\) in each entry. Extracting it from every row
will give a cubic saving in the logarithm, provided the remaining
moments with denominator \(F_K\) are small enough.

\begin{lemma}\label{lem:gamma}
Define \(\Gamma(r)=L_s((1-rs)^{-1})\), with \(L\) acting on the
variable \(s\). This function is holomorphic on \(|r|<|w|^{-1}\),
and for \(1\le a\le K\),
\[
                   L^\sigma_{-\Xi}\big((z-r_a)^{-1}\big)=\Gamma(r_a).
\]
\end{lemma}
\begin{proof}
Use \((z-r_a)^{-1}=-r_a^{-1}(1-w^{-a}z)^{-1}\). By
\eqref{eq:conj} the conjugate of the right side of
\eqref{eq:poles}, with \(\Xi\) replaced by the indeterminate \(X\),
is \(2\tau^a(X+S_a)\). At \(X=-\Xi\) this gives
\[
 L^\sigma_{-\Xi}\big((z-r_a)^{-1}\big)
   =\frac{2\tau^a}{r_a}(\Xi-S_a)
   =2\rho^{-a}(\Xi-S_a)
   =L_s\big((1-r_a s)^{-1}\big).
\]
The middle equality uses \(r_a=(\rho\tau)^a\). Direct differentiation
in the original functional gives
\[
                 \Gamma(r)=2r\sum_{n\ge1}
                       \frac{\tau^{3n}}{(1-rw^n)^3},
\]
locally uniformly on the indicated disk.
\end{proof}

Take \(K=2\) and \(j=0\), so the rational function is \(1/F_2\),
with poles \(r_1=w\) and \(r_2=w^2\). By Lemma~\ref{lem:gamma} the
two simple fractions contribute
\[
  \frac{\Gamma(r_1)}{r_1-r_2}+\frac{\Gamma(r_2)}{r_2-r_1}
     =\frac{\Gamma(r_1)-\Gamma(r_2)}{r_1-r_2}.
\]
For a fixed holomorphic \(\Gamma\), this divided difference tends to
\(\Gamma'(r_1)\) as \(r_2\to r_1\). Estimating the two fractions
separately loses that cancellation. For \(K\) nodes, a single contour
will combine all the pole contributions as the nodes accumulate at zero.

\begin{proposition}\label{prop:second}
At the prescribed value of the conjugate variable,
\[
          \log|\Delta_{K,\sigma\tau}(-\Xi)|
             \le -K^2(K+1)\log\eta+O(K\log K).
\]
\end{proposition}
\begin{proof}
Write \(Y_j=L^\sigma_{-\Xi}(z^j/F_K(z))\) for \(0\le j\le 2K-2\).
By Lemma~\ref{lem:gamma} and the residue theorem, the contribution
from the simple fractions is the counterclockwise contour integral
\begin{equation}\label{eq:contour}
 \sum_{a=1}^K\frac{r_a^j\Gamma(r_a)}{F'_K(r_a)}
       =\frac{1}{2\pi i}\int_{|z|=1}\frac{z^j\Gamma(z)}{F_K(z)}\,dz .
\end{equation}
On this circle \(|F_K(z)|\ge\prod_{a\ge1}(1-|w|^a)>0\), and
\(\Gamma\) is bounded there. Thus \eqref{eq:contour} is bounded
independently of \(j,K\). This uses the signed analogue of the
common pole function in \cite[Proposition~5.1]{KLucas}.

It remains to bound the polynomial quotient when \(j\ge K\).
On expanding \(1/F_K(z)\) at infinity its coefficients after the
leading \(z^{-K}\) are \(h_\ell=h_\ell(r_1,\ldots,r_K)\), defined by
\(\sum_{\ell\ge0}h_\ell t^\ell=\prod_{a=1}^K(1-r_a t)^{-1}\).
The polynomial quotient is
\(\sum_{\ell=0}^{j-K}h_\ell z^{j-K-\ell}\), and
\(\sum_{\ell\ge0}|h_\ell|\le\prod_{a\ge1}(1-|w|^a)^{-1}\).
Also \(\sigma\alpha_i=Q/\alpha_i\), so
\(\sigma\nu_i=Q i(i+1)/(\alpha_i-Q)\) is \(O((i+1)^2)\),
using \(|\alpha_i|\le u<1\). Consequently the entire moment
\(Y_j\) is \(O((K+1)^2)\).

Returning to \eqref{eq:denconj} multiplies the absolute determinant
of the \(Y_{i+j}\) by \(|w|^{K^2(K+1)/2}\). Hadamard's inequality
and \(|w|=\eta^{-2}\) give the result.
\end{proof}

\section{Nonvanishing at arbitrarily large sizes}\label{sec:nonzero}

The Cauchy transform of \(L\) is the function from
Lemma~\ref{lem:gamma} in a reciprocal variable. For \(|z|>|w|\),
\[
  \mathcal C(z)=L_s\big((z-s)^{-1}\big)
       =\frac{1}{z}L_s\big((1-s/z)^{-1}\big)
       =\frac{\Gamma(1/z)}{z}.
\]
The points \(r_a=w^a\) of Section~\ref{sec:second} correspond to
\(x_a=w^{-a}\). These exterior points are the zeros of \(D_K\)
for \(1\le a\le K\), and \(|x_a|\to\infty\).

Kronecker proved that the generating function of a fixed moment sequence
is rational exactly when its Hankel determinants eventually vanish
\cite[Arts.~VII and~IX]{Kronecker}. Zudilin recalls this criterion in
the irrationality setting \cite[Section~5]{Zudilin}. Here the moments
change with \(K\), but every zero of \(D_K\) is also a zero of \(D_{K+1}\).
We use these shared points in a
rational interpolation argument, following \cite[Proposition~5.1]{QZeta}
and \cite[Proposition~6.1]{KLucas}, for the actual signed functional.

\begin{proposition}\label{prop:nonzero}
For arbitrarily large \(K\), one has \(\Delta_{K,\tau}(\Xi)\ne0\).
\end{proposition}
\begin{proof}
The series for \(\Gamma\) in Lemma~\ref{lem:gamma} gives the continuation
\[
 \mathcal C(z)=\sum_{n\ge1}\frac{2\rho^n w^n z}{(z-w^n)^3}.
\]
The series converges normally on compact subsets away from the nodes
and their accumulation point \(0\). Its cubic pole coefficient at
\(w^n\) is \(2\rho^n w^{2n}\ne0\). It follows that \(\mathcal C\)
has infinitely many poles and is not rational.

Suppose now that all determinants from some size onward vanish.
At each such size \(k\), take a nonzero polynomial \(R_k\) of
degree \(<k\) in the kernel. Thus
\(L_s(s^jR_k(s)/D_k(s))=0\) for \(0\le j<k\).
We seek a polynomial \(T_k\) such that the error
\(\mathcal C R_k-T_k\) vanishes at \(x_1,\ldots,x_k\) and is
\(O(1/z)\) at infinity. Consecutive sizes share the first \(k\)
points. With this decay those shared zeros will force
\(T_k/R_k=T_{k+1}/R_{k+1}\). To obtain both properties, set
\begin{equation}\label{eq:interp}
 T_k(z)=L_s\!\left(
   \frac{R_k(z)D_k(s)-D_k(z)R_k(s)}{(z-s)D_k(s)}
                      \right),\qquad
 E_k(z)=\mathcal C(z)R_k(z)-T_k(z).
\end{equation}
The numerator in \eqref{eq:interp} vanishes at \(z=s\), so \(T_k\)
is a polynomial of degree \(<k\). Subtracting in \eqref{eq:interp}
gives
\[
 E_k(z)=D_k(z)L_s\!\left(\frac{R_k(s)}{(z-s)D_k(s)}\right).
\]
Thus \(E_k\) vanishes at \(x_a\), \(1\le a\le k\),
where the functional on the right is analytic. At infinity expand
\((z-s)^{-1}\) in powers of \(1/z\). The expansion converges with
two \(s\)-derivatives on a disk containing the nodes, and the first
\(k\) terms give zero by the kernel equations. The factor \(D_k(z)\)
has degree \(k\), so \(E_k(z)=O(1/z)\).

The polynomial
\[
          T_{k+1}R_k-T_kR_{k+1}=R_{k+1}E_k-R_kE_{k+1}
\]
has degree at most \(k-1\) by its growth at infinity, and has
the \(k\) common zeros \(x_1,\ldots,x_k\). Hence it vanishes
identically. All sufficiently large \(k\) therefore give the same
rational function \(T_k/R_k=T/R\), written in lowest terms.
Write \(R_k=V_kR\), \(T_k=V_kT\) with \(V_k\) a nonzero polynomial.
At a zero of \(V_k\), cancelling this common factor can lose an
interpolation condition. The decay at infinity bounds how many can
be lost. Indeed, the fixed error \(E=\mathcal C R-T\) is nonzero since
\(\mathcal C\) is not rational. Write \(E(z)\sim c z^{-q}\) at
infinity, with \(c\ne0\). As \(E_k=V_kE=O(1/z)\), we have
\(\deg V_k\le q-1\). But \(E(x_a)\ne0\) for all sufficiently large
\(a\), since \(|x_a|\to\infty\). Every such \(x_a\) with \(a\le k\)
must be a zero of \(V_k\). Their number eventually exceeds the fixed
bound \(q-1\), a contradiction.
\end{proof}

\newpage
\section{The norm and the fixed-stride sums}

\begin{proof}[Proof of Theorem~\ref{thm:main}]
Suppose \(A^-_{P,Q}=v/m\) with \(v\in\Z,\ m\ge1\).
By \eqref{eq:rational}, \(mD^2\Xi\in\Ord\). The degree bound after
\eqref{eq:det} and the coefficient clearing in \eqref{eq:clearing}
show that
\[
                y_K=(mD^2)^K C_K(\tau)\Delta_{K,\tau}(\Xi)\in\Ord .
\]
No cyclotomic factor vanishes at \(\tau\), since \(|\tau|<1\).
By Proposition~\ref{prop:nonzero} we can choose arbitrarily large
\(K\) with \(y_K\ne0\). For these indices \(\Nm_{E/\Q}(y_K)\)
is a nonzero integer. Under our rationality assumption
\(\sigma\Xi=-\Xi\) by \eqref{eq:rational}, so its absolute value is
\[
 |\Nm_{E/\Q}(y_K)|=(mD^2)^{2K}
 |C_K(\tau)C_K(\sigma\tau)|\,
 |\Delta_{K,\tau}(\Xi)|\,|\Delta_{K,\sigma\tau}(-\Xi)|.
\]
The first factor contributes only \(2K\log(mD^2)\) to the logarithm.
For the other three factors use Lemma~\ref{lem:normcost} and
Propositions~\ref{prop:first} and~\ref{prop:second}, respectively. This gives
\[
 0\le \log|\Nm_{E/\Q}(y_K)|
 \le -\frac{296}{11025}K^3\log\eta
                  +O(K^2\log K)+2K\log(mD^2),
\]
by \eqref{eq:margin}. The right-hand side tends to \(-\infty\)
because \(\eta>1\) and \(m,D\) are fixed. This proves irrationality.
\end{proof}

\begin{proof}[Proof of Corollary~\ref{cor:stride}]
Write \(\bar\eta=Q/\eta\) and \(P_k=\eta^k+\bar\eta^k\). This is an
integer, by the recurrence for the traces of the powers of \(\eta\).
Binet's formula gives
\[
                    U_{kn}(P,Q)=U_k(P,Q)\,U_n(P_k,Q^k).
\]
If \(Q^k=1\), then \(P_k=\eta^k+\eta^{-k}>2\). If \(Q^k=-1\),
then \(P_k=\eta^k-\eta^{-k}>0\). Thus the new pair is admissible.
Apply Theorem~\ref{thm:main} and divide by the nonzero integer \(U_k^3\).
For Fibonacci, \(U_2(1,-1)=1\) and \(P_2=3\), so \(F_{2n}=U_n(3,1)\).
This gives the irrational sum \(F_2^{-3}-F_4^{-3}+F_6^{-3}-\cdots\).
\end{proof}

\section*{Verification}
Another AI model independently audited the mathematics and exposition of
the 18 November 2026 working proof, including the fixed-stride result.
Two independent reviews subsequently examined the exposition of the
seven-page manuscript of the same date.

\clearpage
\begin{thebibliography}{99}
\bibitem{DNNS}
D.~Duverney, Ke.~Nishioka, Ku.~Nishioka and I.~Shiokawa,
Transcendence of Rogers--Ramanujan continued fraction and reciprocal
sums of Fibonacci numbers, \emph{Proc.\ Japan Acad.\ Ser.~A Math.\
Sci.} \textbf{73} (1997), 140--142.
\href{https://doi.org/10.3792/pjaa.73.140}{doi:10.3792/pjaa.73.140}.
\bibitem{Kyoto}
D.~Duverney, Ke.~Nishioka, Ku.~Nishioka and I.~Shiokawa,
Transcendence of Rogers--Ramanujan continued fraction and reciprocal
sums of Fibonacci numbers, \emph{RIMS K\^oky\^uroku} \textbf{1060}
(1998), 91--100. \href{https://www.kurims.kyoto-u.ac.jp/~kyodo/kokyuroku/contents/pdf/1060-10.pdf\#page=5}{Theorem~4, p.~95}.
\bibitem{ESS}
C.~Elsner, S.~Shimomura and I.~Shiokawa, Algebraic relations for
reciprocal sums of Fibonacci numbers, \emph{Acta Arith.}
\textbf{130} (2007), 37--60.
\href{https://doi.org/10.4064/aa130-1-3}{doi:10.4064/aa130-1-3}.
\bibitem{KFib}
S.~Krisman, The sum of the reciprocal cubes of the Fibonacci numbers
is irrational, \href{https://hexagonmath.org/2611.00473v1}{Hexagon 2611.00473v1} (2026).
\bibitem{KLucas}
S.~Krisman, The reciprocal cube sums of Lucas sequences of quadratic
units are irrational, \href{https://hexagonmath.org/2611.00491v1}{Hexagon 2611.00491v1} (2026).
\bibitem{Zudilin}
W.~Zudilin, A determinantal approach to irrationality,
\emph{Constr.\ Approx.} \textbf{45} (2017), 301--310.
\href{https://doi.org/10.1007/s00365-016-9333-7}{doi:10.1007/s00365-016-9333-7}.
\bibitem{Uvarov}
V.~B.~Uvarov, The connection between systems of polynomials orthogonal
with respect to different distribution functions,
\emph{USSR Comput.\ Math.\ Math.\ Phys.} \textbf{9} (1969), no.~6, 25--36.
\href{https://doi.org/10.1016/0041-5553(69)90124-4}{doi:10.1016/0041-5553(69)90124-4}.
\bibitem{Kratt}
C.~Krattenthaler, A determinant identity for moments of orthogonal
polynomials that implies Uvarov's formula for the orthogonal
polynomials of rationally related densities,
\href{https://arxiv.org/abs/2103.03969}{arXiv:2103.03969} (2021).
\bibitem{Kronecker}
L.~Kronecker, Zur Theorie der Elimination einer Variabeln aus zwei
algebraischen Gleichungen, \emph{Monatsber.\ K\"onigl.\ Preuss.\
Akad.\ Wiss.\ Berlin} (1881), 535--600.
\href{https://www.e-rara.ch/download/pdf/5761162.pdf\#page=42}{Art.~IX in the collected works}.
\bibitem{QZeta}
uplift-universal-inc, Irrationality of the Lambert-series value
\(\zeta_{1/p}(3)\), research draft (2026).
\href{https://github.com/uplift-universal-inc/qzeta_3-irrationality}{Manuscript and accompanying sources}.
\bibitem{checks}
twinheterodox, Verification programs accompanying \emph{The alternating sum of
reciprocal Fibonacci cubes is irrational} (2026).
Included with the source as electronic supplementary material.
\href{https://github.com/twinheterodox/alternating-lucas-cubes}{Project repository address}.
\end{thebibliography}
\end{document}
